% Section 10 -- Conclusion
\section{Conclusion}
\label{sec:conclusion}

We have presented \textbf{BGP Tracer}, an Agentic
two-phase analyser for the joint attribution of
Internet routing and traffic anomalies. The headline
result on the curated ground-truth (29 events, 3
classes) is \textbf{26/29 = 89.7\%} strict Top-1
attribution, \textbf{Macro-F1 0.866}, and per-cause
recall of 92.9\% (outage), 87.5\% (leak), and
85.7\% (hijack), with the principled arbitration
contributing \textbf{55.2\,pp} over Pick-Max (34.5\%)
and \textbf{16.4\,pp} over a gradient-boosted
ML-oracle (73.3\%). The five-month live deployment
confirms production-rate throughput; the released
artifact generates every figure in this paper from raw
telemetry in a single shell command, making the
paper's results fully reproducible.

The broader contribution is not the headline accuracy but
the demonstration that causal attribution in production BGP
networks is tractable without a known causal graph --- and that
the reason the classical toolkit appears to fail on this data is
a question-mismatch, not a defect in the tests. Granger, PC, and
transfer entropy return refusals on BGP telemetry because their
preconditions are systematically violated by bursty,
non-stationary update streams; they answer ``does the past of
routing linearly improve prediction of traffic,'' which is not
what an operator asks. BGP Tracer instead anchors attribution on
a pre/post difference-in-means estimator whose intervention
window is the routing breakpoint, and routes the classical tests
through a Bayesian posterior where disagreement costs a channel
its weight rather than its vote. The abstention discipline is
what makes the approximation auditable: the engine commits only
when the estimator's significance conditions hold and the
synthesized posterior clears the commit floor, and it records
the triggering statistic whenever it does not, so a refusal
is a verdict with a stated reason.

We are explicit about where the causal evidence actually comes
from. On the 29 live-stream events the Pearl-style estimator
itself satisfies the commit gate on only 3; the remaining 26 are
resolved by the routing-withdrawal metadata signal and the
domain-informed prior. That is a real dependency on domain
knowledge, not a purely statistical result, and we state it
rather than let a 29/29 headline imply otherwise. Closing that
gap --- either by discovering a BGP causal graph that supplies an
adjustment set, or by widening traffic coverage so the pre/post
contrast is estimable on more events --- is the most valuable
next step.

Future work: a held-out population back-filled to 2022--2024
(the current skew follows Cloudflare Radar's retention window),
multi-collector deployment to remove the single-vantage
assumption, and integration of a discovered causal graph for
tighter identification bounds.
